\documentclass[11pt,letterpaper]{article}

\usepackage[margin=0.82in]{geometry}
\usepackage[T1]{fontenc}
\usepackage{lmodern}
\usepackage{microtype}
\usepackage{xcolor}
\usepackage{booktabs}
\usepackage{tabularx}
\usepackage{longtable}
\usepackage{enumitem}
\usepackage{listings}
\usepackage{hyperref}
\usepackage{fancyhdr}
\usepackage{titlesec}

\definecolor{ink}{HTML}{172033}
\definecolor{blue}{HTML}{245C8A}
\definecolor{pale}{HTML}{EEF4F8}
\definecolor{warning}{HTML}{8A4B20}
\hypersetup{colorlinks=true,linkcolor=blue,urlcolor=blue,citecolor=blue,
  pdftitle={Goal Relevance Governance for Persistent AI Work},
  pdfauthor={ZeroBot with Benjamin Goertzel}}

\pagestyle{fancy}
\fancyhf{}
\lhead{Goal Relevance Governor}
\rhead{Design note -- 2026-08-14}
\cfoot{\thepage}
\setlength{\headheight}{14pt}
\setlist{nosep,leftmargin=1.5em}
\titleformat{\section}{\Large\bfseries\color{ink}}{\thesection}{0.6em}{}
\titleformat{\subsection}{\large\bfseries\color{blue}}{\thesubsection}{0.6em}{}

\lstdefinestyle{plain}{
  basicstyle=\ttfamily\small,
  backgroundcolor=\color{pale},
  frame=single,
  rulecolor=\color{blue!30},
  breaklines=true,
  columns=fullflexible,
  keepspaces=true
}

\newcommand{\term}[1]{\textbf{#1}}
\newcommand{\verdict}[1]{\texttt{#1}}

\title{\vspace{-1.4cm}\textbf{Goal Relevance Governance for Persistent AI Work}\\
\large From textual LLM task graphs to Atomspace and PLN reasoning}
\author{ZeroBot, in discussion with Benjamin Goertzel}
\date{August 14, 2026}

\begin{document}
\maketitle

\begin{abstract}
Persistent AI workers can be locally competent yet globally irrational. They
may continue producing commits, tests, and reports even after their parent goal
has been achieved, superseded, or demoted; they may impose a safe but very slow
execution cadence without exposing the speed--risk tradeoff; or they may block
urgent work because resource conflicts are invisible outside the local task.
This note proposes a \emph{Goal Relevance Governor}: a typed, evidence-bearing
graph connecting top-level goals, intermediate goals, tasks, resources,
results, and constraints. Bounded reviewers evaluate these links before work,
after material events, and periodically. Each review must yield an operational
verdict such as continue, accelerate, pause, stop, replan, or escalate.

The graph is augmented by a stage-sensitive \emph{result contract}. It states
the project kind, current maturity stage, next decision-relevant result,
evidence required to trust that result, intended mature result, current
non-goals, and conditions under which hardening becomes necessary. This
prevents exploratory research from silently turning into production
engineering and prevents generalized infrastructure from displacing a simpler
path to an urgent operational result.

The initial implementation can use committed JSON or Markdown plus LLM
critics. A later incarnation can represent the same abstract interface in an
Atomspace (a typed knowledge graph) and use symbolic languages plus
probabilistic reasoning. OmegaHive Deliberation Rooms, a proposed format for
bounded multi-agent discussion, remain complementary: they provide the social surface for
ambiguous strategic review, while the committed graph provides authoritative
state and detects when deliberation is needed. The proposal is motivated by
concrete failures in recent OpenClaw work and includes a shadow-first
validation plan based on retrospective replay. Product and project names are
defined when they first appear; understanding the proposal does not require
prior knowledge of the systems that motivated it.
\end{abstract}

\tableofcontents
\newpage

\section{The problem: competent locally, stupid globally}

Current agent infrastructure is relatively good at checking local execution:
did a command run, did tests pass, was an artifact written, is a worker alive,
and did a promised deliverable land? These checks are necessary. They do not
answer a more important question:

\begin{quote}
Is this task, in its present form, still a good means to a currently important
goal, given what is now known and what else it delays or blocks?
\end{quote}

A task can generate abundant \term{activity evidence} while generating no
\term{parent-goal progress}. This distinction is load-bearing. Commits, passing
tests, completed subtasks, model calls, and status reports prove activity.
Goal progress requires evidence that the desired world state has become closer,
less uncertain, less risky, or less costly to reach.

Persistent workers amplify the gap. A well-scoped prompt, once scheduled,
tends to preserve its local objective. Without a means--ends review layer, the
worker can faithfully execute an obsolete or strategically dominated plan for
days. The scheduler sees health; the human experiences drift.

\subsection{The missing feedback loop}

The missing loop has four components:

\begin{enumerate}
  \item maintain explicit top-level goals and priorities;
  \item connect tasks to those goals through reusable intermediate goals;
  \item update the links when results, assumptions, costs, or priorities change;
  \item repeatedly test whether active tasks remain relevant and well-shaped.
\end{enumerate}

This is not full general intelligence. It is a disciplined approximation to
cross-task means--ends reasoning, opportunity-cost awareness, and reflective
control. It can remove a consequential class of failures even when each local
worker remains an ordinary LLM agent.

\subsection{How to read the project-specific examples}

The motivating incidents occurred in a research workspace containing
persistent software agents, scheduled experiments, and several related
symbolic-AI tools. The names are incidental to the argument:

\begin{itemize}
  \item \term{OpenClaw} is the host agent platform that schedules and runs
        persistent workers.
  \item \term{OmegaClaw} is an experimental family of conversational research
        agents built from an upstream codebase plus preserved memories and
        selected capabilities.
  \item \term{PeTTa} and \term{SWI-Prolog} are parts of the symbolic-language
        runtime used by some experiments and agents.
  \item \term{Chroma} is a vector database used to store and retrieve agent
        memories.
\end{itemize}

Readers can mentally substitute ``scheduled implementation worker,''
``long-running research experiment,'' ``shared runtime,'' ``production
service,'' and ``memory database.'' The proposed governor is not specific to
any of these products.

\section{Motivation from observed history}

The examples below are not hypothetical reconstructions. They are concrete
failures that motivated the proposal. They also show why task correctness is
not enough: each task had a locally intelligible rationale.

\subsection{Migration fragmented into ten-minute slices}

A migration intended to finish quickly was decomposed into short, repeated
execution slices. Short slices offered recoverability, frequent observation,
and bounded damage. But they also imposed repeated startup, rehydration,
inspection, reporting, and scheduling overhead. The worker optimized a local
safety mechanism without periodically comparing it against the parent goal:
complete the migration promptly.

The failure was not simply choosing safety over speed. The deeper failure was
not presenting the tradeoff. A relevance review should have compared at least
three alternatives:

\begin{itemize}
  \item frequent short slices: highest observability, highest overhead;
  \item one continuous run with checkpoints: lower overhead and recoverable;
  \item one uninterrupted run: fastest, with the largest recovery cost.
\end{itemize}

Given a high priority on completion speed, the system should have proposed the
continuous checkpointed path without waiting for the human to notice the
cadence problem.

\subsection{An implementation worker pursuing an achieved milestone}

A scheduled code-generation worker, internally called Plain2MeTTa, continued
targeting an earlier demonstration milestone that had already been completed
and immutably tagged. The worker
continued to emit nominally successful runs and encountered incidental shell
and test-path failures, but no current top-level goal was advanced. The actual
next goal was a revised second-version intermediate representation for the
translation pipeline.

This should be a nearly deterministic stop condition. If a task's sole parent
goal is marked achieved, the task is unschedulable until it is closed or linked
to a new unmet goal. A simple graph query would have exposed the stale edge:

\begin{lstlisting}[style=plain]
task: scheduled-code-generation-worker
  contributes_to -> milestone:initial-demonstration [status=achieved]
  contributes_to -> (none active)
verdict: STOP_STALE; request retargeting
\end{lstlisting}

\subsection{A long-running research experiment blocking urgent service restoration}

A long-horizon computational-chemistry experiment, internally called
petta-chem, occupied a shared symbolic-language runtime and source tree. A
clean restoration of three broken conversational agents was treated as blocked
until that experiment finished, even though service restoration was urgent
and the chemistry run was checkpointed and resumable. Once the conflict was
made explicit, the experiment was paused after 242 of 276 durable work units,
its completed results were preserved, and the agent-restoration work could
proceed.

The missing information was not buried in either task. It was the cross-task
relation:

\begin{lstlisting}[style=plain]
task: checkpointed-chemistry-calibration
  occupies -> resource:shared-symbolic-runtime
  supports -> goal:long-horizon-research

task: clean-agent-service-restoration
  requires -> resource:shared-symbolic-runtime-quiescent
  supports -> goal:restore-conversational-agents [priority=urgent]

derived: chemistry-calibration blocks service-restoration
recommended verdict: PAUSE_RECOVERABLY(chemistry-calibration)
\end{lstlisting}

The correct response did not require deep chemistry or runtime intelligence.
It required representing resource contention, priority, and reversibility in
one decision context.

\subsection{An over-engineered agent-repair path}

The agent-repair effort itself accumulated a privileged process inspector,
several layers of launch wrappers, synthetic safety-boundary tests, and
repeated local review gates. Many components addressed real safety concerns.
Yet the chosen path became dominated by machinery whose completion was not
necessary for the top-level goal: restore three usable conversational agents
from a clean copy of the maintained upstream software while preserving
valuable memories.

A relevance review should repeatedly ask whether a simpler path now dominates:
fresh upstream installation, three isolated agent identities, ordinary
conversation tests, disposable copies of the memory database for compatibility
testing, and selective migration of useful capabilities.
This illustrates \term{plan-level staleness}: every subtask may remain useful
in isolation while the overall decomposition becomes an inferior route.

\subsection{Premature hardening of research infrastructure}

Several research and agent-infrastructure lanes accumulated increasingly
specific hardening tasks: path and encoding guards, lifecycle edge cases,
bounded audit behavior, generalized supervisors, and ever larger regression
suites. Many changes were individually sound. The alignment question was
often left implicit: was the hardened component already known to support a
scientifically or operationally valuable result?

For exploratory research, the next useful result is usually epistemic. It may
be a calibrated estimator, a planted-signal recovery, a meaningful baseline
comparison, a falsifying counterexample, or a first credible scientific
signal. Generalized production hardening has weak value if the underlying
method may soon be discarded. The correct early standard is not carelessness;
it is \term{minimum sufficient rigor}: enough correctness, provenance,
controls, reproducibility, cost containment, security, and recoverability to
make the next result trustworthy.

This distinction challenges hardening tasks with a counterfactual:

\begin{quote}
If this hardening is omitted now, does the next scientific or operational
result become invalid, uninterpretable, irreproducible, unsafe, unaffordable,
or unrecoverable?
\end{quote}

If not, the task should normally be deferred until evidence justifies a
reusable research tool or production system. A large safety suite can be
essential for a production delegation substrate while being irrelevant to an
earlier question such as whether a proposed algorithm produces any useful
signal.

\section{Other failure classes the governor should catch}

The same structure applies beyond the observed incidents.

\begin{longtable}{p{0.24\textwidth}p{0.68\textwidth}}
\toprule
Failure class & Example and expected detection \\
\midrule
\endhead
Goal completion drift & Tests keep improving a feature after its acceptance
criterion passed. Detect active tasks whose parent goals are achieved. \\
Goal supersession & A new architecture replaces the old one, but old workers
continue. Detect edges to superseded goals and plans. \\
Priority inversion & A low-value benchmark monopolizes a GPU, database, or
human review channel required by an urgent incident. Derive a blocks relation. \\
Proxy gaming & A worker increases test counts or commits while the actual
quality metric is unchanged. Separate activity from goal-delta evidence. \\
Sunk-cost continuation & Repeated failures accumulate, but prior investment is
treated as a reason to continue. Recompute prospective value, not historical cost. \\
Risk drift & A once-safe task becomes dangerous after scope, permissions, or
external state changes. Trigger review on assumption and resource changes. \\
Cost drift & API usage, runtime, or human attention exceeds the initial
estimate. Trigger review on material estimate variance. \\
Dependency drift & An upstream release obviates custom work, or a dependency
becomes unavailable. Re-evaluate necessary-for and alternative-path edges. \\
Premature optimization & Performance work begins before correctness or
estimator validation. Enforce prerequisite and necessary-for relations. \\
Coordination overload & Excessive checkpoints, deliberations, or reports slow
the goal they protect. Model governance itself as a cost-bearing task. \\
Human bottleneck creation & The agent requests repeated copy/paste, approval,
or archaeology that could be automated. Represent human attention as scarce. \\
Silent tradeoff choice & The worker selects safety, speed, quality, or cost
without reporting the alternatives. Require explicit tradeoff evidence. \\
\bottomrule
\end{longtable}

\section{Stage-sensitive goals and result contracts}

Top-level priorities are necessary but do not completely specify rational
work. A goal such as ``reach a mature useful result quickly'' can be
interpreted as permission to polish indefinitely. The governor therefore
needs three general meta-goals plus an explicit definition of the result being
pursued now.

\subsection{Three sharpened meta-goals}

\begin{description}[style=nextline]
  \item[A: Decision-relevant speed]
    Reach the next result that can rationally change belief, direction, or
    action as quickly as possible.
  \item[B: Least commitment]
    Prefer the simplest, most reversible, and least-committing path that can
    produce a trustworthy result.
  \item[C: Mature-result efficiency]
    Reach the intended mature result as quickly as feasible, avoiding shortcuts
    whose probability-weighted downstream cost exceeds their present value.
\end{description}

B is an Occam heuristic, not an absolute command. ``Simplest'' means simplest
subject to the evidential and operational validity required at the current
stage. It does not justify a shortcut that makes an experiment uninterpretable
or risks destruction of valuable data. C prevents A from becoming naive
short-termism, but C must not be used to fund speculative maturity before the
project has earned it.

\subsection{Useful results depend on project kind}

The phrase \emph{useful result} has different meanings in different projects.

\begin{longtable}{p{0.22\textwidth}p{0.31\textwidth}p{0.39\textwidth}}
\toprule
Project kind & Typical next useful result & Typical mature useful result \\
\midrule
\endhead
Exploratory research & A decision-changing signal, falsification, calibrated
probe, or meaningful baseline result & A replicated and well-delimited
scientific conclusion; hardened software may remain a separate project \\
Confirmatory research & Independent replication, stronger controls, or a
predeclared test & A robust result with uncertainty, ablations, limitations,
and reproducible artifacts \\
Prototype engineering & A thin end-to-end slice proving the central user or
system capability & A reusable implementation adequate for repeated trials \\
Production engineering & A bounded service path under representative load & A
secure, observable, maintainable, fault-tolerant deployed system \\
Service recovery & Restoration of the smallest standard known-good path & A
stable restored service with preserved state and selectively reintroduced
features \\
Infrastructure research & Evidence that the substrate enables a valuable
downstream task & A reusable substrate whose hardening is proportional to its
demonstrated use \\
\bottomrule
\end{longtable}

An exploratory project's mature result is often a successful or negative
scientific conclusion, not hardened code. Once the scientific value is clear,
building a reusable or production implementation can become a distinct project
with its own goal, budget, acceptance criteria, and risks.

\subsection{A maturity ladder}

A useful default ladder is:

\begin{description}[style=nextline]
  \item[E0 -- Exploration]
    Run the smallest trustworthy experiment that can falsify or support the
    core idea.
  \item[E1 -- Signal validation]
    Calibrate the estimator, add decisive controls, and establish one credible
    result.
  \item[E2 -- Replication and robustness]
    Test seeds, datasets, baselines, ablations, perturbations, and uncertainty.
  \item[E3 -- Reusable research implementation]
    Add the cleanup, documentation, interfaces, and tests required for repeated
    scientific use.
  \item[E4 -- Production engineering]
    Add security, fault tolerance, deployment, monitoring, performance, and
    broad edge-case hardening.
\end{description}

The labels are defaults, not a universal ontology. A service-recovery project
may use a parallel ladder: unchanged known-good baseline, isolated instances,
state compatibility, selective feature return, and production soak. What
matters is that each active task names the stage it advances.

\subsection{The result contract}

Each active project or major goal should maintain:

\begin{enumerate}
  \item project kind;
  \item current phase or maturity stage;
  \item next decision-relevant result;
  \item evidence sufficient for that result to be trustworthy;
  \item intended mature result;
  \item known shortcuts and their likely debt;
  \item conditions under which additional hardening becomes necessary;
  \item explicit non-goals for the current phase;
  \item owner, provenance, confidence, and next review trigger.
\end{enumerate}

The contract is not frozen forever. A meaningful result, failure, cost change,
or new alternative should trigger reconsideration. An LLM can propose the
contract from project records, but consequential definitions and priority
changes should initially be confirmed by the human owner.

\subsection{Minimum-sufficient-hardening test}

Hardening at E0 or E1 is presumed deferred unless omitting it materially
threatens one of the following for the next result:

\begin{itemize}
  \item scientific validity or estimator correctness;
  \item interpretability of the result;
  \item reproducibility and provenance;
  \item protection against silent data corruption;
  \item cost containment or recoverability of expensive work;
  \item security, privacy, or explicit authorization boundaries;
  \item the ability to run the discriminating experiment at all.
\end{itemize}

If the test passes, implement the smallest guard that protects the result. If
it fails, record a deferred-hardening edge and the event that should reactivate
it. This preserves useful technical debt as an explicit option rather than
either ignoring it or paying it prematurely.

\subsection{Comparing the first result with the mature result}

The governor should compare a small set of alternative plans using at least:

\begin{itemize}
  \item expected time to the next decision-relevant result;
  \item probability that the result is valid and informative;
  \item probability-weighted rework and technical debt;
  \item risk of losing data, violating constraints, or creating lock-in;
  \item opportunity cost imposed on other goals;
  \item governance and coordination overhead;
  \item option value retained by a reversible prototype.
\end{itemize}

The estimates may initially be qualitative. Their purpose is to expose the
tradeoff, not manufacture false precision. A continuous checkpointed migration
may dominate both tiny repeated slices and an unsafe uninterrupted run. A
small amount of architecture can reduce total time to E3, but only when its
expected value exceeds the chance that E0 shows the idea should be abandoned.

\subsection{Applying the contract to conversational-agent recovery}

For the failed service-repair effort, the result contract should have read:

\begin{lstlisting}[style=plain]
project_kind: service_recovery
current_stage: unchanged_known_good_software_baseline
next_result:
  ordinary bounded multi-turn conversation on a pinned upstream version
trust_evidence:
  deterministic smoke, bounded lifecycle, no orphan processes
mature_result:
  three reliable isolated agents with preserved memories and selected skills
current_non_goals:
  generalized process inspector; custom deferred runner; broad supervisor
\end{lstlisting}

Here ``pinned upstream'' means an exact, recorded version of the maintained
base software, before custom extensions are restored. Against that contract,
a generalized privileged process inspector contributes mainly
to a hypothetical later operational maturity stage. A clean pinned upstream
installation tests the central uncertainty directly, is more reversible, uses
fewer assumptions, and reaches the next result sooner. The correct verdict is
\verdict{DEFER} for the inspector and \verdict{DO\_FIRST} for the clean
baseline.

\section{The proposed control structure}

\subsection{A typed, evidence-bearing graph}

The minimal persistent object is a directed property graph. Text can be the
initial semantic payload, but node and edge types should be explicit.

\paragraph{Core node types.}
\begin{itemize}
  \item \term{Goal}: a desired world state with priority, provenance, time
        horizon, success criteria, status, and authority.
  \item \term{Intermediate goal}: a reusable enabling state between strategic
        goals and concrete work.
  \item \term{Task}: an executable commitment with owner, expected duration,
        cost, risk, reversibility, acceptance test, and evidence path.
  \item \term{Result}: an observed outcome, artifact, measurement, failure, or
        decision that updates beliefs about the graph.
  \item \term{Resource}: CPU, GPU, process tree, repository, identity,
        credential, database, network endpoint, or human attention.
  \item \term{Constraint}: safety, budget, authorization, deadline, ordering,
        privacy, semantic, or lifecycle invariant.
  \item \term{Plan}: a candidate decomposition or path through intermediate
        goals and tasks, allowing competing plans to be compared.
  \item \term{Result contract}: project kind, active stage, next result, trust
        evidence, mature result, current non-goals, and hardening trigger.
\end{itemize}

\paragraph{Core edge types.}
\begin{itemize}
  \item \texttt{contributes-to}, \texttt{necessary-for}, and
        \texttt{sufficient-for};
  \item \texttt{depends-on}, \texttt{blocks}, \texttt{conflicts-with}, and
        \texttt{competes-for};
  \item \texttt{achieves}, \texttt{invalidates}, \texttt{supersedes}, and
        \texttt{provides-evidence-for};
  \item \texttt{alternative-to}, \texttt{refines}, and \texttt{part-of};
  \item \texttt{authorized-by}, \texttt{proposed-by}, and
        \texttt{reviewed-by}.
\end{itemize}

Each uncertain edge should carry confidence, provenance, assumptions, last
review time, and a criterion for revisiting it. A bare weight is inadequate:
two equal weights can mean very different things if one rests on direct
evidence and the other on an LLM guess.

\subsection{Priorities are not merely scalar}

A weighted score is useful for ranking, but some relationships should be
lexicographic or constraint-based. For example:

\begin{enumerate}
  \item do not violate explicit safety and authorization constraints;
  \item address urgent production restoration before discretionary research;
  \item among allowed tasks, maximize expected progress per scarce resource;
  \item prefer simpler, more reversible plans when expected value is similar.
\end{enumerate}

An illustrative task value is

\[
V(t) = P(g)\,U(t,g)\,S(t)\,T(t)
       - C(t) - R(t) - O(t) - G(t),
\]

where $P(g)$ is inherited goal priority, $U(t,g)$ expected contribution,
$S(t)$ probability of success, and $T(t)$ urgency. $C(t)$ is direct cost,
$R(t)$ risk, $O(t)$ opportunity cost imposed on other paths, and $G(t)$
governance/coordination overhead. This is a prompt for structured judgment,
not initially a claim of cardinal precision.

\subsection{Required review questions}

Every active task should periodically answer:

\begin{enumerate}
  \item Which active goal does success advance, through which intermediate path?
  \item What new parent-goal delta occurred since the previous review?
  \item Is the parent goal achieved, superseded, paused, or reprioritized?
  \item Is a simpler, faster, cheaper, or safer route now available?
  \item What resources are held, and which higher-priority paths are delayed?
  \item Which assumptions or estimates have changed?
  \item What are the principal speed--safety--quality--cost tradeoffs?
  \item Which result contract and maturity stage does this task advance?
  \item Does proposed hardening pass the minimum-sufficient-hardening test?
  \item Should the task continue, accelerate, defer, pause, stop, replan, or escalate?
\end{enumerate}

\subsection{Operational verdicts}

A review that only writes analysis will not change behavior. It must emit one
of a small number of verdicts with reasons, evidence, expiry, and authority:

\begin{description}[style=nextline]
  \item[\verdict{CONTINUE}] Current plan remains proportionate and relevant.
  \item[\verdict{ACCELERATE}] Reduce cadence/governance overhead or allocate
        more resources within authorization.
  \item[\verdict{PAUSE\_RECOVERABLY}] Preserve state and release a resource.
  \item[\verdict{STOP\_STALE}] Parent goal is achieved, superseded, or absent.
  \item[\verdict{REPLAN}] A better path or invalidated assumption changes the decomposition.
  \item[\verdict{DEFER}] The task may matter at a later maturity stage but is
        not necessary for the next decision-relevant result.
  \item[\verdict{ESCALATE}] Human or multi-seat judgment is required.
  \item[\verdict{BLOCKED}] No safe authorized step is available; name the blocker.
\end{description}

\section{When relevance is reviewed}

Periodic review is necessary but insufficient. The system should review on
events that can change the means--ends relation:

\begin{itemize}
  \item before a worker starts or resumes;
  \item after a meaningful success, failure, or acceptance result;
  \item when a goal is created, achieved, paused, superseded, or reprioritized;
  \item when the next result, project kind, or maturity stage changes;
  \item when expected duration, cost, or risk changes materially;
  \item when a new resource conflict or blocking relation appears;
  \item after repeated activity with no recorded parent-goal delta;
  \item at a bounded periodic cadence as a backstop.
\end{itemize}

The cadence should itself be adaptive. A stable low-risk task with valid
assumptions needs fewer reviews than an incident response with rapidly changing
state. Reviews should be triggered by semantic events where possible, not by
arbitrary ten-minute fragmentation.

\section{Relation to OmegaHive Deliberation Rooms}

Deliberation Rooms and the Goal Relevance Governor solve different layers of
the problem.

\begin{center}
\begin{tabularx}{\textwidth}{>{\bfseries}p{0.24\textwidth}XX}
\toprule
Dimension & Goal Relevance Governor & Deliberation Rooms \\
\midrule
Primary question & What work remains a good means to current goals? & How do
humans and seats reason together and land a conclusion? \\
Durable state & Goal/task/resource graph and verdict evidence & Open/close
events, committed conclusion references, archived transcript \\
Execution shape & Cheap bounded evaluator, event-triggered or periodic &
Bounded wake-based participation with chair and synthesizer \\
Authority & Initially advisory; later policy-gated actions & Conversation has
no authority until conclusions are committed \\
Failure detected & Stale tasks, priority inversion, blocked resources,
inferior plans & Unlanded deliberation, attribution and synthesis failures \\
\bottomrule
\end{tabularx}
\end{center}

The integration is:

\begin{lstlisting}[style=plain]
committed goal/task/resource graph
              |
              v
bounded relevance evaluator
       |                    |
 clear low-risk verdict     ambiguous / novel / high-impact
       |                    |
       v                    v
commit recommendation     open Deliberation Room
                              |
                     seats + human deliberate
                              |
                     synthesizer commits result
                              |
                              v
                    graph is re-evaluated
\end{lstlisting}

Most checks should not open rooms. Making every review conversational would
recreate coordination overhead. A room is appropriate when uncertainty,
conflict, novelty, impact, or authority exceeds a threshold. Examples include
changing a top-level priority, abandoning a major research direction, choosing
between materially different safety profiles, or resolving disagreement among
project seats.

The graph lives in the committed spine. A room discusses proposed changes;
only the synthesizer's committed artifact changes authoritative goal state.
This follows the Deliberation Rooms rule that speech is not real until it lands
through the write path.

\section{A simple LLM-centric OpenClaw incarnation}

\subsection{Storage and projections}

The first version should be intentionally small:

\begin{itemize}
  \item committed JSON or YAML records for goals, nodes, edges, and verdicts;
  \item stable identifiers and append-only event history;
  \item projections generated from existing \texttt{PROJECT.md},
        \texttt{TASKS.md}, \texttt{DECISIONS.md}, cron jobs, sessions, and
        resource observations;
  \item textual descriptions on nodes, with typed metadata around them;
  \item source pointers for every factual status and uncertain inference.
\end{itemize}

Project Markdown remains the human-readable source of project facts. The graph
can initially be a normalized cross-project projection rather than a second
competing truth store. Explicit goal and priority records are added only where
the current project files lack them.

\subsection{LLM roles}

Separate LLM calls can perform bounded functions:

\begin{enumerate}
  \item \term{Link proposer}: suggests intermediate goals and typed edges.
  \item \term{Result-contract proposer}: classifies project kind and stage,
        proposes next and mature results, sufficient evidence, and non-goals.
  \item \term{Evidence checker}: verifies whether claims have source support.
  \item \term{Relevance critic}: answers the required review questions.
  \item \term{Alternative planner}: proposes simpler or faster paths.
  \item \term{Adversarial reviewer}: searches for priority inversion,
        opportunity cost, hidden blockers, and silent tradeoffs.
  \item \term{Synthesizer}: emits a structured recommendation or prepares a
        Deliberation Room packet.
\end{enumerate}

These may initially be prompts to one strong model, but their contracts should
remain distinct. The evaluator must not silently edit human priorities. It may
propose changes with confidence and provenance.

\subsection{Authority ladder}

\begin{itemize}
  \item \term{Read-only shadow}: record what would have been recommended.
  \item \term{Advisory}: notify the operator of stale or conflicting work.
  \item \term{Scheduling gate}: prevent a worker with no active goal path from
        starting, while allowing explicit override.
  \item \term{Low-risk autonomy}: pause checkpointable work or retarget an
        obviously stale worker under predeclared rules.
  \item \term{Strategic autonomy}: defer until extensive validation; major goal
        and priority changes remain human-authoritative.
\end{itemize}

This progression prevents a relevance governor from becoming another source
of global damage. It should earn authority through counterfactual evidence.

\section{A symbolic reasoning incarnation}

\subsection{Atomspace representation}

An Atomspace is a typed knowledge graph whose nodes and links can themselves
participate in reasoning. The same interface can be represented as atoms
rather than JSON objects. A schematic MeTTa-like symbolic vocabulary might
include:

\begin{lstlisting}[style=plain]
(Goal RestoreConversationalAgents)
(Priority RestoreConversationalAgents urgent)
(Status RestoreConversationalAgents active)

(Task ChemistryCalibration)
(Task CleanAgentInstallation)
(ContributesTo ChemistryCalibration LongHorizonResearch)
(ContributesTo CleanAgentInstallation RestoreConversationalAgents)
(Occupies ChemistryCalibration SharedSymbolicRuntime)
(Requires CleanAgentInstallation QuiescentSharedSymbolicRuntime)
(Checkpointable ChemistryCalibration True)

(ProjectKind AgentRecovery ServiceRecovery)
(CurrentStage AgentRecovery KnownGoodBaseline)
(NextResult AgentRecovery OrdinaryBoundedConversation)
(MatureResult AgentRecovery ThreeIsolatedAgentsWithPreservedMemory)
(CurrentNonGoal AgentRecovery GeneralizedProcessInspector)

(Evidence E1 (CompletedUnits ChemistryCalibration 242 276))
(Confidence (Blocks ChemistryCalibration CleanAgentInstallation) 0.99)
\end{lstlisting}

Intermediate goals become first-class atoms. Rules can derive transitive
contribution, blocked paths, alternative paths, and goal achievement. Temporal
atoms can represent deadlines, expiry, last review, and changing status.

\subsection{PLN-style uncertain inference}

Probabilistic Logic Networks (PLN), a framework for uncertain logical
inference over such graphs, can add more than a symbolic transcription. It can combine uncertain
means--ends links and evidence:

\begin{itemize}
  \item infer the strength with which a task contributes to a goal through
        several uncertain intermediate links;
  \item weaken a contribution when repeated results fail to change the parent
        state;
  \item compare competing plans using uncertain success and cost estimates;
  \item reason separately about time-to-information and time-to-mature-result;
  \item propagate new evidence that invalidates assumptions;
  \item distinguish lack of evidence from evidence of irrelevance;
  \item explain derived verdicts as inference trails.
\end{itemize}

A simple conceptual rule is:

\begin{quote}
If task $T$ contributes to intermediate goal $I$, and $I$ contributes to goal
$G$, then $T$ contributes to $G$ with an uncertainty derived from both links;
new evidence about either link updates the composite relevance.
\end{quote}

The numerical semantics require care. Contribution is not ordinary truth, and
priority is not probability. It may be better to use several typed values:
belief in the causal link, expected utility, urgency, cost, and confidence in
the estimate. PLN should reason over these without collapsing them into one
opaque scalar too early.

\subsection{Existing components and their roles}

The governor should reuse rather than replace existing components:

\begin{itemize}
  \item \term{ThreadKeeper}, a persistent-worker lifecycle service, supplies
        durable task events and worker-status primitives.
  \item \term{GoalChainer}, an experimental symbolic goal-appraisal module,
        supplies norm/goal/task relations and uncertain belief grading seams.
  \item \term{Petta-memory}, a symbolic memory layer, can preserve evidence and
        attribution needed by uncertain relevance inference.
  \item \term{Deliberation Rooms}, a bounded multi-agent discussion format,
        supply strategic review when routine evaluation is insufficient.
  \item \term{Conversation Governor}, a message-control layer, handles message admission and egress;
        it is adjacent but does not decide which tasks are strategically useful.
\end{itemize}

The abstraction seam is the typed graph plus verdict contract. Storage,
inference, and transport can change without rewriting scheduling policy.

\section{Validation before authority}

\subsection{Retrospective replay corpus}

The initial detector should be evaluated against frozen historical episodes,
including at least:

\begin{enumerate}
  \item migration slice overhead before the continuous-run proposal;
  \item the stale code-generation worker pursuing an achieved milestone;
  \item a checkpointed chemistry experiment blocking urgent service recovery;
  \item the over-engineered conversational-agent repair path;
  \item premature hardening tasks before evidence of research value;
  \item matched controls where long-running or cautious work was justified.
\end{enumerate}

For each episode, freeze only information available at the review time. A
successful system must flag the problem before the human intervention, name
the correct parent-goal issue, and propose a proportionate verdict. It must not
use later outcomes as hidden evidence.

\subsection{Metrics}

\begin{itemize}
  \item detection recall for known alignment failures;
  \item false-alarm rate on justified long-running tasks;
  \item decision latency from decisive evidence to recommendation;
  \item estimated time and resource savings under counterfactual replay;
  \item human agreement with the stated tradeoff and verdict;
  \item correct classification of project kind, stage, and useful result;
  \item hardening precision: necessary guards retained, premature polish deferred;
  \item explanation completeness and evidence traceability;
  \item governor overhead relative to savings;
  \item rate of reviews producing no new information.
\end{itemize}

The last two protect against a recursive failure: building a large governance
system whose cost exceeds the drift it prevents.

\subsection{Shadow-first deployment}

Deployment should proceed from offline replay to read-only shadow observation,
then advisory recommendations, and only later to narrow scheduling gates. No
live autonomy is justified merely because the architecture is persuasive.
Predeclared promotion thresholds and explicit rollback should govern each
authority increase.

\section{Design risks and safeguards}

\subsection{Bureaucratic overgrowth}

The graph can become a ceremonial ontology maintained more carefully than the
work itself. Start with the few node/edge types needed to detect actual
failures. Every field should support a query or decision. Measure governance
overhead.

The same warning applies to maturity models. The E0--E4 ladder is a diagnostic
aid, not a reason to create five phases of ceremony for every small task.

\subsection{LLM confabulation and unstable links}

LLMs may invent causal relationships or rewrite priorities in rhetorically
plausible ways. Separate proposals from accepted links; attach evidence,
confidence, author, and expiry; use adversarial review for major edges; and
keep human priorities authoritative at first.

\subsection{Priority laundering}

Derived scores can conceal value judgments. Preserve the original human goal
and priority statement, show every inherited priority path, and expose the
tradeoff dimensions before aggregation.

\subsection{Over-centralization}

A global governor can become a brittle bottleneck. Projects should retain
local autonomy within explicit resource and priority bounds. The global layer
intervenes mainly on cross-project conflicts, stale parent goals, and major
plan changes.

\subsection{Premature enforcement}

Incorrect pauses may destroy momentum or invalidate experiments. Begin
read-only; prefer recoverable pauses; require explicit authority for
irreversible actions; and retain human override with a recorded rationale.

\section{A staged implementation path}

\subsection*{Phase 0: Design and frozen examples}
Complete this specification; define a small retrospective corpus and expected
verdicts; freeze result contracts for the motivating cases; identify existing
data sources and invariants.

\subsection*{Phase 1: Read-only graph projection}
Project current goals, tasks, cron workers, sessions, and resource conflicts
into a typed JSON graph. Include project kind, maturity stage, next result, and
current non-goals. Do not alter scheduling.

\subsection*{Phase 2: LLM relevance critic}
Run event-triggered and periodic reviews. Emit structured shadow verdicts with
evidence and alternatives. Evaluate replay metrics.

\subsection*{Phase 3: Deliberation integration}
Open rooms for ambiguous or high-impact cases. Require committed synthesis and
measure landing rate and decision latency.

\subsection*{Phase 4: Narrow scheduling gates}
With explicit approval, block only deterministic stale cases, such as a task
whose sole parent goal is achieved. Preserve overrides and rollback.

\subsection*{Phase 5: Atomspace/PLN realization}
Map the stable graph and verdict interface into atoms in a symbolic knowledge
graph. Compare symbolic and LLM judgments on the same corpus. Use PLN for uncertain
contribution and evidence propagation where it demonstrably improves results.

\section{Open research questions}

\begin{enumerate}
  \item What is the smallest graph vocabulary that catches the motivating
        failures without becoming a general project-management ontology?
  \item Can an LLM reliably infer a project's kind, current stage, and next
        decision-relevant result from incomplete project records?
  \item How should contribution, belief, utility, priority, urgency, and risk
        be represented without inappropriate scalar collapse?
  \item How can intermediate goals be reused without erasing project-specific
        semantics?
  \item What event thresholds trigger review without reproducing excessive
        checkpoint overhead?
  \item How should a relevance evaluator estimate opportunity cost when
        alternative plans are only partially known?
  \item How should technical debt and option value be estimated without
        systematically favoring either premature hardening or reckless prototypes?
  \item Which verdicts can safely execute automatically, and which should
        always open a Deliberation Room or request human approval?
  \item Can embedding-based conceptual similarity propose useful reusable
        intermediate goals while provenance and logic prevent false analogy?
  \item Does PLN improve calibration and explanation over a strong LLM critic,
        or mainly provide a more inspectable representation?
\end{enumerate}

\section{Conclusion}

The motivating failures share a single structure: local tasks continued under
fixed assumptions while the global meaning of the work changed. A Goal
Relevance Governor addresses this by making means--ends links, priorities,
resources, evidence, and alternatives explicit and repeatedly reviewable. Its
essential output is not another report but an operational verdict tied to
current goals.

The result contract adds a necessary temporal and epistemic dimension. It
distinguishes the result needed now from a hypothetical mature artifact, and
it distinguishes validity-critical rigor from premature hardening. The
governor can therefore optimize for fast trustworthy learning without losing
sight of downstream maturity.

The practical first version does not require a mature symbolic AGI. Committed
textual records, a small typed graph, bounded LLM critics, and retrospective
validation can already catch stale goals, priority inversions, blocked
resources, and inferior decompositions. The same architecture can later move
into Atomspace and PLN without changing its outer contract.

Deliberation Rooms strengthen rather than replace the proposal. They provide
a bounded venue where humans and seats exercise consequential judgment. The
graph and review workers decide when that judgment is needed and preserve its
committed consequences. Together they form a plausible incremental path from
task-local LLM competence toward persistent reflective intelligence across
projects.

\enlargethispage{4\baselineskip}
\section*{Source and provenance notes}

This document synthesizes Ben Goertzel's 2026-08-14 proposal and the immediate
discussion with ZeroBot. The Deliberation Rooms comparison is based on Cassio
Pennachin, \emph{OmegaHive Deliberation Rooms: Vision, a Slack MVP, and the
relation to the Conversation Governor}, July 29, 2026, preserved locally with
SHA-256
\texttt{8c6d6b3916d7353b5e7818fab2a20abe98d8c8a6d553df936316fec44419e762}.
Historical examples are operational observations from the shared OpenClaw
workspace; they are used here as design evidence, not as claims that the
proposed mechanism has already been validated.

\end{document}
